\section{Conclusions}

This work presents a conservative DBS formulation for PHREEQC-coupled reactive
transport with cell-based surface complexation in image-based porous media,
validated across analytical, image-based, and multi-scale porous-media cases.

The key contribution is a cell-based representation of unresolved reactive
surface area that sits inside the DBS micro-continuum framework. One fixed
porosity field serves aqueous accumulation, transport, and the PHREEQC cell
volumes alike, which keeps the aqueous transport conservative and the local
mass balance closed cell by cell. The unresolved micropore surface enters as an
additive intrinsic volumetric density, so reactive uptake can be represented in
image phases whose internal surface area lies inside the coarse cell rather
than on its boundary.

The 1D tests preserve the closed-batch Ca inventory to a relative residual
below $10^{-9}$ and show approximately first-order mesh and timestep
convergence, while the surface-area reconstruction recovers analytic interface
area to sub-percent accuracy for 3D spheres at $R/\Delta x\ge16$ and
quasi-2D circular posts at all four tested resolutions. In the Li micromodel, the
intrinsic-area representation keeps both aqueous and sorbed inventories within
$0.6\%$ of the resolved result over $Pe=0.1$--10, while the sorbed streamwise
profile differs by up to $4.75\%$. Across the 2D and 3D image-based cases, the
model resolves the dissolved and adsorbed fields that follow from the local
pore geometry and competition between species for surface sites.

Surface complexation can therefore be carried conservatively on cells whose
reactive area is unresolved, freeing the reactive surface from the resolved
mesh boundary in multi-scale image-based simulation. This matters for reactive
porous-media problems in which selectivity and surface accessibility control
performance, among them contaminant retention, radionuclide migration, and
critical-mineral recovery by sorbent-bearing porous media. Constraining the intrinsic surface-area input against measured internal area is
the step that would most improve predictive accuracy, and extending the
formulation to evolving reactive area and kinetic surface chemistry would open
the dissolution and precipitation problems that the equilibrium treatment
excludes.
